\subsection{Eigenvalues and Eigenvectors}

% **Role:** structure revealed by invariant subspaces and eigenbehavior.
%
% Eigenvalues, eigenvectors, eigenspaces, invariant subspaces, triangularization, diagonalization.

Spectral linear algebra understands an operator through subspaces it preserves. Eigenstructure, triangularization, diagonalization, and commuting families all refine this invariant-subspace viewpoint.

\begin{definition} \label{def:eigenvalue-eigenvector} \hlt{Eigenvalue and Eigenvector}\\
Let $T\in\LL(\mathcal{V})$. A scalar $\lambda\in\mathbb{F}$ is an \hlt{eigenvalue} of $T$ if there exists $\mathbf{v}\in\mathcal{V}$ with $\mathbf{v}\neq\mathbf{0}$ and $T\mathbf{v}=\lambda\mathbf{v}$. Such a vector $\mathbf{v}$ is an \hlt{eigenvector} of $T$ corresponding to $\lambda$.
\end{definition}

\begin{definition} \label{def:eigenspace} \hlt{Eigenspace}\\
Let $T\in\LL(\mathcal{V})$ and $\lambda\in\mathbb{F}$. The \hlt{eigenspace} of $T$ corresponding to $\lambda$ is $E_{\lambda}(T)\defeq\ker(T-\lambda I_{\mathcal{V}})$. Thus $\lambda$ is an eigenvalue of $T$ if and only if $E_{\lambda}(T)\neq\{\mathbf{0}\}$. Equivalently, $E_{\lambda}(T)$ is the nullspace of the shifted operator $T-\lambda I_{\mathcal{V}}$.
\end{definition}

\begin{definition} \label{def:spectrum-and-geometric-multiplicity} \hlt{Spectrum and Geometric Multiplicity}\\
Let $T\in\LL(\mathcal{V})$. The \hlt{spectrum} of $T$ is $\operatorname{spec}(T)\defeq\{\lambda\in\mathbb{F}:\lambda\text{ is an eigenvalue of }T\}$.\\
The \hlt{geometric multiplicity} of $\lambda\in\mathbb{F}$ is $\dim E_{\lambda}(T)$. For a square matrix $A\in\operatorname{Mat}(n,n,\mathbb{F})$, $E_{\lambda}(A)\defeq\ker(A-\lambda I_n)$ and $\operatorname{spec}(A)$ denotes the spectrum of the induced operator on $\mathbb{F}^n$.
\end{definition}

\begin{proposition} \label{prop:eigenvectors-as-kernel-vectors} \hlt{Eigenvectors as Kernel Vectors}\\
Let $T\in\LL(\mathcal{V})$, $\lambda\in\mathbb{F}$. Vector $\mathbf{v}$ is eigenvector of $T$ corresponding to $\lambda$ if and only if $\mathbf{v}\in\ker(T-\lambda I_{\mathcal{V}})$, $\mathbf{v}\neq\mathbf{0}$.
\begin{proof}
\begin{enumerate}[label=\roman*.]
\setlength{\itemsep}{0pt}
\item[$(\Rightarrow)$] If $\mathbf{v}$ is an eigenvector corresponding to $\lambda$, then $\mathbf{v}\neq\mathbf{0}$ and $T\mathbf{v}=\lambda\mathbf{v}$, so $(T-\lambda I_{\mathcal{V}})\mathbf{v}=\mathbf{0}$.
\item[$(\Leftarrow)$] If $\mathbf{v}\in\ker(T-\lambda I_{\mathcal{V}})$ and $\mathbf{v}\neq\mathbf{0}$, then $(T-\lambda I_{\mathcal{V}})\mathbf{v}=\mathbf{0}$, so $T\mathbf{v}=\lambda\mathbf{v}$.
\end{enumerate}
\hlt{Structural Commentary}: Eigenvectors are kernels of shifted operators.\\
\hlt{Key Insight}: Subtract $\lambda I_{\mathcal{V}}$ to convert $T\mathbf{v}=\lambda\mathbf{v}$ into homogeneous form.
\end{proof}
\end{proposition}

\begin{theorem} \label{thm:equivalent-eigenvalue-conditions} \hlt{Equivalent Conditions for an Eigenvalue}\\
Let $\mathcal{V}$ be finite-dimensional, let $T\in\LL(\mathcal{V})$, and let $\lambda\in\mathbb{F}$. The following are equivalent:
\begin{enumerate}[label=(\roman*)]
\setlength{\itemsep}{0pt}
\item $\lambda$ is an eigenvalue of $T$;
\item $T-\lambda I_{\mathcal{V}}$ is not injective;
\item $T-\lambda I_{\mathcal{V}}$ is not surjective;
\item $T-\lambda I_{\mathcal{V}}$ is not invertible.
\end{enumerate}
\begin{proof}
\begin{enumerate}[label=(\roman*).]
\setlength{\itemsep}{0pt}
\item[$(i)$]$\Rightarrow (ii)$ If $\lambda$ is an eigenvalue, Proposition \ref{prop:eigenvectors-as-kernel-vectors} gives nonzero vector in $\ker(T-\lambda I_{\mathcal{V}})$, so $T-\lambda I_{\mathcal{V}}$ not injective.
\item[$(ii)$]$\Rightarrow (i)$ If $T-\lambda I_{\mathcal{V}}$ not injective, then $\ker(T-\lambda I_{\mathcal{V}})$ contains nonzero vector, so $\lambda$ is eigenvalue (Prop \ref{prop:eigenvectors-as-kernel-vectors}).
\item[$(ii)$]$\Rightarrow (iii)$ If $T-\lambda I_{\mathcal{V}}$ were surjective, Theorem \ref{thm:finite-dimensional-operator-invertible-equivalences} would make it injective, a contradiction.
\item[$(iii)$]$\Rightarrow (iv)$ If $T-\lambda I_{\mathcal{V}}$ were invertible, then it would be surjective, a contradiction.
\item[$(iv)$]$\Rightarrow (ii)$ If $T-\lambda I_{\mathcal{V}}$ were injective, Theorem \ref{thm:finite-dimensional-operator-invertible-equivalences} would make it invertible, a contradiction.
\end{enumerate}
\hlt{Structural Commentary}: In finite dimensions, noninvertibility has equivalent forms.\\
\hlt{Key Insight}: A scalar $\lambda$ is spectral precisely when $T-\lambda I_{\mathcal{V}}$ has a nontrivial kernel.
\end{proof}
\end{theorem}

\begin{theorem} \label{thm:distinct-eigenvalues-independent-eigenvectors} \hlt{Eigenvectors for Distinct Eigenvalues Are Independent}\\
Let $T\in\LL(\mathcal{V})$. Suppose $\lambda_1,\ldots,\lambda_m$ are distinct eigenvalues of $T$, and let $\mathbf{v}_1,\ldots,\mathbf{v}_m$ be corresponding eigenvectors. Then $(\mathbf{v}_1,\ldots,\mathbf{v}_m)$ is linearly independent.
\begin{proof}
Suppose the list is linearly dependent. Let $k$ be smallest positive integer s.t. $\mathbf{v}_k\in\operatorname{span}(\mathbf{v}_1,\ldots,\mathbf{v}_{k-1})$.\\
Then $k\geq2$, $(\mathbf{v}_1,\ldots,\mathbf{v}_{k-1})$ is linearly independent, and $\exists$ scalars $a_1,\ldots,a_{k-1}$ s.t. $\mathbf{v}_k=a_1\mathbf{v}_1+\cdots+a_{k-1}\mathbf{v}_{k-1}$.\\
Applying $T$ gives $\lambda_k\mathbf{v}_k=a_1\lambda_1\mathbf{v}_1+\cdots+a_{k-1}\lambda_{k-1}\mathbf{v}_{k-1}$.
Multiplying first displayed equation by $\lambda_k$ and subtracting gives $0=a_1(\lambda_k-\lambda_1)\mathbf{v}_1+\cdots+a_{k-1}(\lambda_k-\lambda_{k-1})\mathbf{v}_{k-1}$.
Since $\mathbf{v}_1,\ldots,\mathbf{v}_{k-1}$ are linearly independent and $\lambda_k\neq\lambda_i$ for $i<k$, all $a_i=0$. Hence $\mathbf{v}_k=\mathbf{0}$, contradicting that $\mathbf{v}_k$ is an eigenvector.\\
\hlt{Structural Commentary}: A smallest dependence isolates the first new eigenvector that depends on earlier ones. Applying $T$ and subtracting separates the distinct eigenvalue factors.\\
\hlt{Key Insight}: Distinct eigenvalues make the same dependence relation impossible after applying $T$.
\end{proof}
\end{theorem}

\begin{corollary} \label{cor:number-of-distinct-eigenvalues} \hlt{Number of Distinct Eigenvalues}\\
If $\mathcal{V}$ is finite-dimensional and $T\in\LL(\mathcal{V})$, then $T$ has at most $\dim(\mathcal{V})$ distinct eigenvalues.
\begin{proof}
If $T$ had more than $\dim(\mathcal{V})$ distinct eigenvalues, choose $m>\dim(\mathcal{V})$ of them, one eigenvector for each.\\
By Theorem \ref{thm:distinct-eigenvalues-independent-eigenvectors}, these eigenvectors form linearly independent list of length $m$, contradicting Theorem \ref{thm:finitedimlenofvec}.\\
\hlt{Structural Commentary}: Distinct eigenvalues give independent directions.\\
\hlt{Key Insight}: Count one eigenvector per eigenspace.
\end{proof}
\end{corollary}

\begin{corollary} \label{cor:invertible-iff-zero-not-eigenvalue} \hlt{Invertibility and the Zero Eigenvalue}\\
Let $\mathcal{V}$ be finite-dimensional and let $T\in\LL(\mathcal{V})$. Then $T$ is invertible if and only if $0\notin\operatorname{spec}(T)$.
\begin{proof}
Apply Theorem \ref{thm:equivalent-eigenvalue-conditions} with $\lambda=0$. Then $0$ is an eigenvalue exactly when $T-0I_{\mathcal{V}}=T$ is not invertible.\\
\hlt{Structural Commentary}: The zero eigenvalue is the spectral form of singularity.\\
\hlt{Key Insight}: Invertibility is equivalent to having trivial kernel, and $E_0(T)=\ker T$.
\end{proof}
\end{corollary}

\begin{corollary} \label{cor:matrix-invertible-iff-zero-not-spectrum} \hlt{Matrix Invertibility and the Spectrum}\\
Let $A\in\operatorname{Mat}(n,n,\mathbb{F})$. Then $A$ is invertible if and only if $0\notin\operatorname{spec}(A)$.
\begin{proof}
\begin{enumerate}[label=\roman*.]
\setlength{\itemsep}{0pt}
\item[$(\Rightarrow)$] If $A$ is invertible, then $\ker(A)=\{\mathbf{0}\}$, so $0$ is not an eigenvalue of $A$.
\item[$(\Leftarrow)$] If $0\notin\operatorname{spec}(A)$, then $\ker(A)=\{\mathbf{0}\}$. Corollary \ref{cor:matrix-rank-nullity} gives full rank, hence $A$ is invertible.
\end{enumerate}
\hlt{Structural Commentary}: Matrix invertibility is the coordinate form of trivial kernel.\\
\hlt{Key Insight}: A zero eigenvalue is exactly a nonzero solution of $A\mathbf{x}=\mathbf{0}$.
\end{proof}
\end{corollary}
